% GENERATED FILE. Edit canonical lessons, then run npm run generate:books.
% canonical-source-hash: fnv1a64:b7dc258c97b8921c
% canonical-lessons: FA-C12-nam, FA-C12-del, FA-C12-dar, FA-C12-ketab

\chapter{Name, Heart, Door, Book}
\label{ch:fa-name-heart-door-book}
\hlchaptermodality{12}

\begin{chapteropening}
\textbf{By the end of this chapter:} \emph{I can say name, heart, door, and book in Persian, and I can tell which of the four is inherited and which is a later Arabic loan.}

The chapter closes on ketâb, an Arabic loan, then runs all four words back to back — nâm, del, dar, ketâb — sorts them by vocabulary layer, and rebuilds the earlier name exchange the reappearance of nâm reopens.
\end{chapteropening}

\section[nâm]{\fa{نام} — name, the word Persian never borrowed}

\label{lesson:FA-C12-nam}



\begin{quote}
Say \textbf{\fa{زبان}}, starting on its letter \textbf{\fa{ز}}. Then say the very first sentence this track ever built, \textbf{\fa{اسمِ} \fa{من} ... \fa{است}}. Both words in that sentence for \textquotedblleft{}name\textquotedblright{} are about to get company.
\end{quote}

\subsection*{You'll want to know first — one word}
\begin{quote}
\textbf{\fa{نام}} — \emph{nâm} — \textbf{name}
\end{quote}

Three letters, all already yours: \textbf{\fa{ن} \fa{ا} \fa{م}} — \emph{n}, long \emph{â}, \emph{m}.

\begin{cousinweb}
\textbf{\fa{نام}} continues Old Persian \textbf{nāma}, from Indo-European *\emph{h\textsubscript{1}nómn}, \textquotedblleft{}name\textquotedblright{} — one of the most widely shared words in the whole family: English kept it as \textbf{name}, Greek as \textbf{ónoma} ($\to$ \textbf{onomastics}), Italian as \emph{nome}, and Sanskrit as \textbf{nāman}. \textbf{\fa{اسم}} \emph{esm}, the word this track taught first, is a separate, later Arabic loan for the identical idea — so Persian now holds two words for \textquotedblleft{}name,\textquotedblright{} inherited and borrowed, sitting side by side the way \textbf{\fa{خوب}} and \textbf{\fa{وقت}} already showed the track's vocabulary layers working.
\end{cousinweb}

\subsection*{Your turn}
\begin{itemize}
\raggedright
  \item \emph{Say it:} \textbf{nâm} — name
  \item \emph{Connect:} \textbf{nâm} $\leftarrow$ *\emph{h\textsubscript{1}nómn} $\to$ English \textbf{name}, Greek \textbf{ónoma}, Sanskrit \textbf{nāman}
  \item \emph{Contrast:} \textbf{\fa{نام}}, inherited, against \textbf{\fa{اسم}}, the Arabic loan the name-exchange lesson taught first
  \item \emph{Retrieve:} \textbf{\fa{زبان}} — tongue, and language, the word that closed the last chapter
\end{itemize}

\begin{tcolorbox}[breakable,colback=teal!4,colframe=teal!35!black,title={Before you move on}]
What does \textbf{\fa{نام}} mean? (\textbf{Name.}) Which other Persian word means the same thing, and which language did it arrive through? (\textbf{\fa{اسم}}, through Arabic.) Which English word shares \fa{نام}'s own root? (\textbf{Name.})

Source: \href{https://en.wiktionary.org/wiki/\%D9\%86\%D8\%A7\%D9\%85}{Wiktionary: \fa{نام}}.
\end{tcolorbox}
\section[del]{\fa{دل} — heart}

\label{lesson:FA-C12-del}



\begin{quote}
Say \textbf{\fa{نام}} (\emph{nâm}), and its English cousin. This lesson's word has an even more famous one.
\end{quote}

\subsection*{You'll want to know first — one word}
\begin{quote}
\textbf{\fa{دل}} — \emph{del} — \textbf{heart}
\end{quote}

Two letters, both already yours: \textbf{\fa{د} \fa{ل}} — \emph{d}, \emph{l}.

\begin{cousinweb}
\textbf{\fa{دل}} continues Middle Persian \textbf{dil}, from the ancient Indo-European word for heart, *\emph{\'{k}\'{\={e}}r} — the same root English wore down to \textbf{heart}, Armenian kept as \textbf{sirt}, and Sanskrit built into \textbf{hṛdaya}. Persian uses \textbf{\fa{دل}} far beyond anatomy, the way English leans on \emph{heart} for courage and feeling rather than only the organ — a later chapter's word for \textquotedblleft{}want\textquotedblright{} will lean on it too.
\end{cousinweb}

\subsection*{Your turn}
\begin{itemize}
\raggedright
  \item \emph{Say it:} \textbf{del} — heart
  \item \emph{Connect:} \textbf{del} $\leftarrow$ *\emph{\'{k}\'{\={e}}r} $\to$ English \textbf{heart}, Sanskrit \textbf{hṛdaya}, Armenian \textbf{sirt}
  \item \emph{Run through:} \textbf{nâm, del} — name, heart
  \item \emph{Retrieve:} \textbf{\fa{دختر}} — daughter, from a few lessons back
\end{itemize}

\begin{tcolorbox}[breakable,colback=teal!4,colframe=teal!35!black,title={Before you move on}]
What does \textbf{\fa{دل}} mean? (\textbf{Heart.}) Which English word is its direct cousin? (\textbf{Heart.}) Name one other language that kept the same root. (\textbf{Sanskrit} \emph{hṛdaya}, or \textbf{Armenian} \emph{sirt}.)

Source: \href{https://en.wiktionary.org/wiki/\%D8\%AF\%D9\%84}{Wiktionary: \fa{دل}}.
\end{tcolorbox}
\section[dar]{\fa{در} — door}

\label{lesson:FA-C12-dar}



\begin{quote}
Say \textbf{\fa{دل}}, and the root it shares with \textbf{heart}. This lesson's word shares its own root with an equally ordinary English word.
\end{quote}

\subsection*{You'll want to know first — one word}
\begin{quote}
\textbf{\fa{در}} — \emph{dar} — \textbf{door}
\end{quote}

Two letters, both already yours: \textbf{\fa{د} \fa{ر}} — \emph{d}, \emph{r}.

\begin{cousinweb}
\textbf{\fa{در}} continues Old Persian \textbf{duvar-}, \textquotedblleft{}door, gate,\textquotedblright{} from Indo-European *\emph{d\textsuperscript{h}wer-}. Few Indo-European words are this well attested across every branch of the family at once: English \textbf{door}, German \textbf{Tür}, Latin \textbf{foris} ($\to$ \textbf{foreign}, \textquotedblleft{}outside the door\textquotedblright{}), Greek \textbf{thýra}, and Sanskrit \textbf{dvāra} all continue the identical ancient word. \textbf{\fa{در}} is Persian's contribution to a cognate set most Indo-European languages still carry.
\end{cousinweb}

\subsection*{Your turn}
\begin{itemize}
\raggedright
  \item \emph{Say it:} \textbf{dar} — door
  \item \emph{Connect:} \textbf{dar} $\leftarrow$ *\emph{d\textsuperscript{h}wer-} $\to$ English \textbf{door}, German \textbf{Tür}, Latin \textbf{foris}, Sanskrit \textbf{dvāra}
  \item \emph{Run through:} \textbf{nâm, del, dar} — name, heart, door
\end{itemize}

\begin{tcolorbox}[breakable,colback=teal!4,colframe=teal!35!black,title={Before you move on}]
What does \textbf{\fa{در}} mean? (\textbf{Door.}) Name three other languages that kept the same root. (\textbf{English} \emph{door}, \textbf{German} \emph{Tür}, \textbf{Sanskrit} \emph{dvāra}, among others.)

Source: \href{https://en.wiktionary.org/wiki/\%D8\%AF\%D8\%B1}{Wiktionary: \fa{در}}.
\end{tcolorbox}
\section[ketâb]{\fa{کتاب} — book, and the name exchange this chapter reopens}

\label{lesson:FA-C12-ketab}



\begin{quote}
Say \textbf{\fa{در}}, and the four languages that share its root. This chapter's last word breaks that pattern completely.
\end{quote}

\subsection*{You'll want to know first — one word}
\begin{quote}
\textbf{\fa{کتاب}} — \emph{ketâb} — \textbf{book}
\end{quote}

Four letters, all already yours: \textbf{\fa{ک} \fa{ت} \fa{ا} \fa{ب}} — \emph{k}, \emph{t}, long \emph{â}, \emph{b}.

\begin{grammarlens}[title={Grammar Lens: a possessed object, the pattern one more time}]
\begin{quote}
\textbf{\fa{کتابِ} \fa{من،} \fa{لطفاً}.} — \emph{ketâb-e man, lotfan.} — \textbf{My book, please.}
\end{quote}

The same ezafe \textbf{-e} from \textbf{\fa{اسمِ} \fa{من}}, \textbf{\fa{کلیدِ} \fa{من}}, and \textbf{\fa{دخترِ} \fa{من}} links a fourth kind of noun to its owner. Four different words, one grammar habit.
\end{grammarlens}

\begin{cousinweb}
\textbf{\fa{کتاب}} is not inherited Persian at all: it is borrowed whole from Arabic \textbf{kitāb}, \textquotedblleft{}book,\textquotedblright{} built on the three-consonant root \textbf{\fa{ك} \fa{ت} \fa{ب}} (\emph{k-t-b}), \textquotedblleft{}to write.\textquotedblright{} Persian's four new words this chapter show its vocabulary layers running in both directions at once — \textbf{\fa{نام}} is inherited where \textbf{\fa{اسم}}, taught first, is Arabic; \textbf{\fa{کتاب}} is Arabic where Persian has no single inherited word of its own for the idea. \textbf{\fa{دل}} and \textbf{\fa{در}} are older still, each carrying an Indo-European root beside a whole set of English cousins.
\end{cousinweb}

\subsection*{Your turn — the chapter, then the name exchange it reopens}
\begin{itemize}
\raggedright
  \item \emph{Say it:} \textbf{ketâb} — book; then \textbf{ketâb-e man, lotfan} — my book, please
  \item \emph{Run through:} \textbf{nâm, del, dar, ketâb} — all four of this chapter's words
  \item \emph{Sort:} which words are inherited, and which one is Arabic
  \item \emph{Rebuild:} \textbf{\fa{شما} / \fa{تو}} — the register choice; \textbf{\fa{اسم} \fa{شما} \fa{چیست؟}} — the question, joining \emph{chi} and \emph{ast}; \textbf{\fa{خوشوقتم}} — the reply that closes an introduction
  \item \emph{Connect:} \textbf{\fa{نام}}, this chapter's own name-word, back to the exchange that first taught \textbf{\fa{اسم}}
  \item \emph{Retrieve:} \textbf{\fa{پا}} — foot, one more body word from a few lessons back
\end{itemize}

\begin{tcolorbox}[breakable,colback=teal!4,colframe=teal!35!black,title={Before you move on}]
Which language did Persian borrow \textbf{\fa{کتاب}} from? (\textbf{Arabic.}) Say \textquotedblleft{}my book, please.\textquotedblright{} (\textbf{Ketâb-e man, lotfan.}) Which of this chapter's four words is inherited Persian, and which is the Arabic loan? (\textbf{\fa{نام،} \fa{دل،} \fa{در}} inherited; \textbf{\fa{کتاب}} Arabic.) Which earlier word for \textquotedblleft{}name\textquotedblright{} does \textbf{\fa{نام}} now stand beside? (\textbf{\fa{اسم}}.)

Source: \href{https://en.wiktionary.org/wiki/\%DA\%A9\%D8\%AA\%D8\%A7\%D8\%A8}{Wiktionary: \fa{کتاب}}.
\end{tcolorbox}
